% ============================================================================
% Artículo (formato HEP): {{TITULO}}
% Autor:    {{AUTOR}}
% Fecha:    {{FECHA}}
%
% Formato inspirado en la plantilla de documentos de la colaboración LHCb.
% Reproduce sus ajustes tipográficos (clase, márgenes, composición, paquetes),
% no la plantilla oficial ni la identidad institucional del CERN. Los campos
% institucionales van vacíos por defecto: se rellenan con --institucion,
% --num-informe y --publicado-en solo si tu documento de verdad los necesita.
%
% Ejemplo publicado con este formato: arXiv:2509.15873
% ============================================================================
\documentclass[12pt, a4paper]{article}

% --- Idioma y codificación ---
\usepackage[utf8]{inputenc}
\usepackage[T1]{fontenc}
\usepackage[spanish]{babel}

% --- Tipografía y composición ---
\usepackage{lmodern}
\usepackage{microtype}
\raggedbottom
\sloppy

% --- Márgenes ---
\usepackage[
  top=1in,
  bottom=1.25in,
  left=1in,
  right=1in
]{geometry}
\columnsep=5mm

% --- Comportamiento de los flotantes ---
\renewcommand{\textfraction}{0.01}
\renewcommand{\floatpagefraction}{0.8}
\renewcommand{\topfraction}{0.9}
\renewcommand{\bottomfraction}{0.9}

% --- Matemáticas ---
\usepackage{amsmath}
\usepackage{amssymb}
\usepackage{amsfonts}
\usepackage{upgreek}   % griegas rectas: obligatorias en notación de partículas

% --- Figuras y tablas ---
\usepackage{graphicx}
\graphicspath{{figuras/}}
\usepackage{float}
\usepackage{caption}
\captionsetup{font=small, labelfont=bf}
\usepackage{subcaption}
\usepackage{booktabs}
\usepackage{array}
\usepackage{multirow}

% --- Notas al pie al fondo de la página ---
\usepackage[bottom, flushmargin, hang, multiple]{footmisc}

% --- Enlaces ---
\usepackage{xcolor}
\usepackage[
  colorlinks=true,
  linkcolor=black,
  citecolor=black,
  urlcolor=blue
]{hyperref}
\usepackage[all]{hypcap}

% --- Bibliografía ---
\usepackage[
  backend=biber,
  style={{ESTILO_CITAS}},
  sorting={{SORTING_CITAS}}
]{biblatex}
\addbibresource{referencias.bib}

{{LINENO_PAQUETE}}

% --- Utilidades ---
\usepackage{ifthen}
\usepackage{xspace}

% ============================================================================
% Macros de partículas y unidades
%
% Conjunto propio y deliberadamente corto: solo lo que se usa a diario.
% El juego completo de símbolos se distribuye con la plantilla oficial de la
% colaboración; no se reproduce aquí. Añade los tuyos debajo.
% ============================================================================
\newcommand{\Bc}{\ensuremath{B_c^+}\xspace}
\newcommand{\jpsi}{\ensuremath{J/\psi}\xspace}
\newcommand{\Dz}{\ensuremath{D^0}\xspace}
\newcommand{\Dsp}{\ensuremath{D_s^+}\xspace}
\newcommand{\Kp}{\ensuremath{K^+}\xspace}
\newcommand{\pip}{\ensuremath{\pi^+}\xspace}
\newcommand{\GeV}{\ensuremath{\mathrm{\,Ge\kern -0.1em V}}\xspace}
\newcommand{\MeV}{\ensuremath{\mathrm{\,Me\kern -0.1em V}}\xspace}
\newcommand{\GeVcc}{\ensuremath{\mathrm{\,Ge\kern -0.1em V\!/}c^2}\xspace}
\newcommand{\MeVcc}{\ensuremath{\mathrm{\,Me\kern -0.1em V\!/}c^2}\xspace}
\newcommand{\ifb}{\ensuremath{\mbox{\,fb}^{-1}}\xspace}

% ============================================================================
% Metadatos
% ============================================================================
\newcommand{\eltitulo}{{{TITULO}}}
\newcommand{\elautor}{{{AUTOR}}}
\newcommand{\lafecha}{{{FECHA}}}
\newcommand{\lainstitucion}{{{INSTITUCION}}}
\newcommand{\elnuminforme}{{{NUM_INFORME}}}
\newcommand{\lapublicacion}{{{PUBLICADO_EN}}}

\hypersetup{
  pdftitle={\eltitulo},
  pdfauthor={\elautor}
}

% ============================================================================
\begin{document}

% ----------------------------------------------------------------------------
% Portada
% ----------------------------------------------------------------------------
\thispagestyle{empty}

\ifthenelse{\equal{\lainstitucion}{}}{}{%
  \begin{center}
    {\large\scshape \lainstitucion}
  \end{center}
  \vspace{2em}
}

\begin{flushright}
  \ifthenelse{\equal{\elnuminforme}{}}{}{\elnuminforme \\}
  \lafecha
\end{flushright}

\vspace{4em}

\begin{center}
  {\LARGE\bfseries \eltitulo}\\[2.5em]
  {\large \elautor}
\end{center}

\vspace{2.5em}

\begin{center}
  \begin{minipage}{0.85\textwidth}
    \begin{center}
      \textbf{Resumen}
    \end{center}
    \vspace{0.5em}
    \noindent
    Un solo párrafo: qué se midió, con qué datos, con qué método, y el
    resultado principal con su incertidumbre. Sin citas y sin referencias
    a secciones del propio documento.
  \end{minipage}
\end{center}

\vfill

\ifthenelse{\equal{\lapublicacion}{}}{}{%
  \begin{center}
    \lapublicacion
  \end{center}
}

\clearpage

{{LINENO_ACTIVAR}}

% ============================================================================
\section{Introducción}
% ============================================================================

Contexto físico, qué se sabe ya y qué hueco viene a llenar este trabajo.
Termina anunciando qué se mide y con qué datos.

% ============================================================================
\section{Datos y selección}
% ============================================================================

Origen de la muestra, luminosidad integrada, criterios de selección y su
justificación.

% ============================================================================
\section{Método}
% ============================================================================

Modelo de ajuste, estimadores y procedimiento. Todo proceso estocástico
declara su semilla.

% Ejemplo de figura:
% \begin{figure}[htbp]
%   \centering
%   \includegraphics[width=0.8\textwidth]{ejemplo.pdf}
%   \caption{Descripción de la figura. Toda masa se nombra por el sistema al
%            que pertenece.}
%   \label{fig:ejemplo}
% \end{figure}

% ============================================================================
\section{Incertidumbres sistemáticas}
% ============================================================================

Catálogo de fuentes, cómo se estima cada una y cómo se combinan.

% Ejemplo de tabla:
% \begin{table}[htbp]
%   \centering
%   \caption{Descripción de la tabla.}
%   \label{tab:ejemplo}
%   \begin{tabular}{lcc}
%     \toprule
%     Fuente & Valor & Unidad \\
%     \midrule
%     Dato 1 & Dato 2 & Dato 3 \\
%     \bottomrule
%   \end{tabular}
% \end{table}

% ============================================================================
\section{Resultados}
% ============================================================================

Valor central, incertidumbre estadística, incertidumbre sistemática, y en ese
orden.

% ============================================================================
\section{Conclusiones}
% ============================================================================

Qué se concluye y qué queda abierto.

% ============================================================================
% Bibliografía
% ============================================================================
\printbibliography

\end{document}
